\subsection{Condições de inserção orbital}
Logo após o lançamento da espaçonave \emph{chaser} ser realizada por um veículo lançador,
% como o \emph{New Glenn} \cite{BlueOrigin_NewGlenn}, \emph{Vulcano} \cite{ULA_VulcanCentaur} ou até o \emph{Falcon 9 Block 5} \cite{SpaceX_Falcon9},
a primeira condição a ser atendida para que ocorra a transferência de órbita, é a inserção em uma órbita estável ao redor da Terra.
% Neste trabalho, a inserção foi considerada ser em uma órbita circular, pois a condição inicial é simplificada e utilizada como referência em missões de RVD/B.

A condição para que a trajetória seja circular é que a força gravitacional forneça a aceleração centrípeta necessária para manter o movimento, portanto, a velocidade de inserção em uma órbita de raio $r$ é dada por
\begin{equation}
    \label{eq:veloc_insercao_orbital}
v_{orb} = \sqrt{\frac{\mu}{r}},
\qquad
\mu = G * M_{Terra}
\end{equation}
onde $\mu$ é o parâmetro gravitacional da Terra, $G$ é a constante de gravitação universal, $M_{Terra}$ é a massa da Terra, incluindo oceanos e atmosfera \cite{IC2021}. 
A equação \eqref{eq:veloc_insercao_orbital} fornece a condição inicial de velocidade que garante a permanência da espaçonave em órbita circular. 
Após estabelecer as condições de inserção orbital, utiliza-se as equações do movimento orbital, derivadas do problema dos dois corpos, detalhadas em \eqref{eq:movimentoTarget}.